% Generated draft: requires mathematical and source review.
\documentclass[11pt]{article}
\usepackage[margin=25mm]{geometry}
\usepackage{fontspec}
\setmainfont{Latin Modern Roman}
\usepackage{amsmath,amssymb,mathtools,unicode-math}
\setmathfont{Latin Modern Math}
\usepackage{enumitem,needspace,xurl,hyperref}
\setlength{\parindent}{0pt}
\setlength{\parskip}{5pt}
\setlength{\emergencystretch}{3em}
\newcommand{\N}{\mathbb{N}}
\newcommand{\Z}{\mathbb{Z}}
\newcommand{\Q}{\mathbb{Q}}
\newcommand{\R}{\mathbb{R}}
\newcommand{\C}{\mathbb{C}}
\newcommand{\F}{\mathbb{F}}
\newcommand{\A}{\mathbb{A}}
\newcommand{\PP}{\mathbb{P}}
\newcommand{\E}{\mathbb{E}}
\newcommand{\eps}{\varepsilon}
\newcommand{\dd}{\,\mathrm{d}}
\begin{document}
\section*{UnsolvedMath ID: 30005274}
\textbf{Source problem code:} catalogue entry
\subsection{Definitions and mathematical statement}
Set \(f(X)=X^3+X\), so that \(d=\deg f=3\). The supplied record does not define \(W_f(a;p)\), including its exponential phase, summation domain, or normalization. Thus a complete mathematical statement cannot be recovered from this record alone.

The following expresses the recorded target conditionally on supplying that definition. Suppose that, for each relevant prime \(p\), it assigns complex values \(W_f(a;p)\) indexed by \(a\in\F_p^\times\). Interpret the extracted fourth moment as
\[
M_f(p)=\frac1p\sum_{a\in\F_p^\times}|W_f(a;p)|^4.
\]
Here the factor \(1/p\) is retained from the original extract; interpreting the sum as running over nonzero residue classes is a reconstruction because its range is incomplete in the supplied original text.

The requested task is to prove, by more elementary methods described in the source as going no further than the Chebotarev Density Theorem, the following alternative:
\[
\lim_{\substack{p\to\infty\\p\ \mathrm{prime}}}M_f(p)=2,
\]
or there exist \(\delta=\delta(d)>0\), a set \(\mathcal P\) of primes, and constants \(C\geq0\) and \(p_0\) such that
\[
\lim_{x\to\infty}
\frac{\#\{p\in\mathcal P:p\leq x\}}{\pi(x)}
\quad\text{exists and is at least }\delta,
\]
\[
M_f(p)\geq3-Cp^{-1/2}
\qquad(p\in\mathcal P,\ p\geq p_0),
\]
where \(\pi(x)\) counts the primes at most \(x\). The last inequality is the usual lower-bound interpretation of the supplied notation \(M_f(p)\geq3+O(p^{-1/2})\); its implied constant must be independent of \(p\), but the record does not specify its permitted dependencies.

This is a request for an elementary proof of a special case of the source's Theorem 2. That theorem, its hypotheses, and the definition of the permitted proof methods are not supplied. The first alternative is supplied by the reconstructed statement but is absent from the fragmentary original quotation. The formulation above therefore remains conditional and requires comparison with the full source before mathematical acceptance.
\subsection{Short English statement}
For \(f(X)=X^3+X\), seek a proof using methods no deeper than Chebotarev density that the recorded fourth moment tends to \(2\), or has lower bound \(3-O(p^{-1/2})\) on a set of primes of positive natural density. The exponential sum \(W_f(a;p)\) and the full theorem hypotheses are missing, so this is a conditional statement-only draft requiring human mathematical review.
\subsection{Sources}
\begin{itemize}
\item Local UnsolvedMath registry, permanent ID \(30005274\), source problem code \texttt{catalogue entry}, title ``Elementary Fourth-Moment Bounds for Cubic Exponential Sums,'' in the Oberwolfach Reports Open Problems collection. No individual proposer is supplied; the proposed year is recorded as 2022.
\item Supplied authoritative source citation: \emph{Analytic Number Theory (2023)}, Oberwolfach Reports, DOI 10.4171/owr/2022/50, pp. 2895--2960. \url{https://doi.org/10.4171/owr/2022/50}. The record attributes the question to workshop open-problem discussions and refers to Theorem 2, whose full statement is not included.
\item Supplied literature source: ``Oberwolfach Report 50,'' \url{https://ems.press/content/serial-article-files/46986}.
\item Literature assessment dated exactly 2026-08-22: the supplied curation labels the problem open and reports that no elementary proof settling the extracted problem was identified. This is supplied dated curation, not an independently verified assessment of the literature or current status.
\item The registry labels the statement \texttt{corrected\_verified} and reports that it was reconstructed and checked against the official report. These are supplied provenance claims; this draft does not independently verify them. The citation uses 2023, while the DOI and recorded proposed year use 2022; these source details are retained as supplied.
\end{itemize}
\end{document}
